\section{Implication discussion}

The results show that immigration has played a decisive role in offsetting the labour-supply losses associated with population ageing, but they also point to the limits of relying on immigration alone. Ageing affects labour supply not only through population structure, but also through behavioural responses, health, hours worked, and exposure to labour-market shocks. The following subsections therefore discuss three implications of the results: the extent to which immigration compensates for ageing, the potential role of healthy ageing as a complementary strategy, and the risks created by future economic and technological disruptions.


\subsection{Compensating for ageing and the volatility trap}


The age-structure effect among Spanish reduced aggregate labour supply by \(2.55\) million FTE workers over 2005--2025. Immigration more than compensated for this loss: the total foreign-born contribution added \(3.04\) million FTE workers, equal to about \(119.4\%\) of the negative Spanish age-structure effect. In short, without immigration, Spanish ageing would have reduced labour supply substantially; with immigration, the aggregate effect becomes positive. Howerver, this migratory buffer exhibits significant pro-cyclical volatility. As evidenced in the 2010--2014 sub-period (Table 3), the foreign-born contribution turned sharply negative (\(-0.54\) million FTE). This reveals a ``volatility trap'': the primary mechanism for offsetting long-term ageing tends to stall or reverse precisely when the economy faces structural stress \parencite{galvez2022cyclicality, funcas2024vulnerability}.

\vspace{0.7em}\par
The role of immigration should also be assessed in light of alternative adjustment mechanisms, especially the behavioural changes induced by ageing. Contrary to migratory buffers, behavioural adjustments among the Spanish-born population proved more resilient to economic cycles. Higher activity rates added \(2.82\) million FTE workers, which, despite a reduction of \(1.60\) million due to declining hours, resulted in a net positive behavioural contribution of \(1.23\) million FTE. Thus, behavioural adjustment helped offset about \(48.1\%\) of the negative age-structure effect among Spanish. While immigration remains the dominant quantitative offset, its reliability as a stabilizer is contingent on economic growth. This highlights that the requirement for migration to balance ageing could be lowered if policy shifted toward supporting the more stable margin of healthy ageing, potentially decoupling labour supply from the boom-bust cycles of migratory flows \parencite{cepr2025limits}.


\subsection{Investment in healthy ageing as a structural stabilizer}

The importance of healthy ageing as a stabilizer is underscored by recent cross-country evidence. Using data for populations aged 50+ in 41 countries from 2000--2022, \textcite{gruss2025healthyAging} find that a decade of cognitive health gains raises labour-force participation by about 20 percentage points, weekly hours worked by around 6 hours, labour productivity by roughly 30\%, and annual labour earnings by about 35\%. In the Spanish context, such investments are increasingly critical not just for the native population, but for a foreign-born population that is ageing more rapidly than in previous decades.

\vspace{0.7em}\par
The transition from the ``exceptional'' migratory wave of the early 2000s \parencite{arango2013exceptional} to a more settled demographic has altered the foreign-born contribution profile in the 2020s \parencite{bde2023ageingParticipation,caixabank2017labourForceAgeing}. Between 2005 and 2025, the share of foreigners aged 55+ increased from \(11.01\%\) to \(23.85\%\), a rise of \(12.84\) percentage points. This demographic maturation has led to a decline in labour-market intensity: the foreign activity rate fell by \(3.33\) percentage points (from \(74.77\%\) to \(71.44\%\)), and average weekly hours dropped by \(2.58\) hours. Without targeted investments in healthy ageing for this settled population, the ``migration buffer'' may become increasingly inefficient, as the volume of new arrivals would need to rise exponentially just to compensate for the declining intensity of the existing, ageing foreign-born stock \parencite{bde2023ageingParticipation, caixabank2017labourForceAgeing}.


\subsection{Disruption on the horizon and the labour supply}

Immigrants are not only ageing more rapidly than the Spanish-born population, but they are also more exposed to labour-market volatility. Foreign workers were disproportionately affected by the 2008 financial crisis: immigrant employment declined, immigration inflows fell by more than 50\% between 2008 and 2013, and Spain experienced net migration outflows \parencite{rodriguez2016labor}. During COVID-19, the unemployment gap between foreign-born and native workers increased from 6.7 to 10 percentage points by late 2020 \parencite{centaur2021,fasani2023,lacaixa2021}. These facts do not contradict the positive demographic contribution of immigration; rather, they show that immigration is more vulnerable to labour-market disruption. This vulnerability also raises a broader question about how ongoing AI disruption may affect labour-market outcomes for immigrants and for the population more generally. At this stage, the evidence remains fragmented and inconclusive \parencite{pizzinelli2023aiExposure}. Some studies find only small effects, or no clear effect, on labour demand in AI-exposed occupations \parencite{hartley2024labor,Ghosh2025AIExposure}. However, because AI has been shown to raise productivity in several work settings \parencite{DellAcqua2024AIProductivity,Noy2023ProductivityAI}, future increases in output may require fewer additional workers than in the past. If so, Spain may not need to rely on immigration to expand labour supply to the same extent as it did over the last two decades.


\section{Conclusion}

Overall, population ageing substantially reduced Spain's potential labour supply, but behavioural adjustment offset much of this loss. Higher participation, especially among Spanish adults and older workers, helped cushion the demographic shock, although declining hours worked limited the size of this compensation. Immigration filled the remaining gap by expanding labour supply mainly through population growth and age composition. However, this contribution depends on migrants' labour-market integration, sectoral exposure, and employment stability \parencite{preston2014effect,fiorio2023migration,naumann2021population}, making it vulnerable to economic and technological disruptions. For the coming decades, complementary strategies may therefore be needed. In particular, investment in healthy ageing appears promising because it can support longer working lives and strengthen labour‑market attachment among older workers in both the Spanish‑born and foreign‑born populations.
